\documentclass[11pt]{article}

\usepackage[margin=1in]{geometry}
\usepackage{graphicx}
\usepackage{amsmath}
\usepackage{siunitx}
\usepackage{caption}
\usepackage{fancyhdr}
\usepackage{float}
\usepackage{hyperref}

\pagestyle{fancy}
\setlength{\headheight}{14pt}
\fancyhf{}
\lhead{Project: \#2 Circuit Design \#1}
\rhead{\thepage}
\renewcommand{\headrulewidth}{0.4pt}

\title{\vspace{-2cm}Linear Circuits 2 \\ Circuit Design \#2 Lab Report}
\author{Group Members: \\ Name: Juvon Mondesir \\ Name: Kenny Gallmon}
\date{}

\begin{document}

\maketitle
\thispagestyle{fancy}

\section*{Objective}

The goal of this project is to use our knowledge of the breadboard and circuit
analysis to produce a physical interpretation of the example circuits that
follow certain constraints. For each part there were different instructions.

\subsection*{Part A}

\begin{figure}[H]
    \centering
    \includegraphics[width=0.85\textwidth]{figures/fig01_part_a.png}
    \caption{Circuit Design Part A}
    \label{fig:part_a}
\end{figure}

For this part, in circuit A, the constraint is $V_{OA} = K_a \cdot V_{IA}$, and
circuit B's was $V_{OB} = K_b \cdot V_{IB}$. $K_a = 0.75$ and $K_b = 0.6$.
While doing this we needed resistors between 1\,k$\Omega$ to 10\,k$\Omega$ and
had to use the fewest components possible.

\subsection*{Part B}

\begin{figure}[H]
    \centering
    \includegraphics[width=0.85\textwidth]{figures/fig02_part_b.png}
    \caption{Circuit Design Part B}
    \label{fig:part_b}
\end{figure}

For the second part of this project, we are now asked to cascade circuit A and
B from the previous portion. Using this, we figure out whether
$V_{OB} = K_a \cdot K_b \cdot V_{IA}$ and explain why it does not. Following
this, we need to redesign one of the two circuits to not only satisfy that
requirement but also the constraints from before, $V_{OA} = K_a \cdot V_{IA}$
and $V_{OB} = K_b \cdot V_{IB}$. Also keeping in mind to use the fewest
components possible and only resistances between 100\,$\Omega$ and
400\,k$\Omega$.

\subsection*{Part C}

\begin{figure}[H]
    \centering
    \includegraphics[width=0.95\textwidth]{figures/fig03_part_c.png}
    \caption{Circuit Design Part C}
    \label{fig:part_c}
\end{figure}

Lastly, for Part C you will use Circuit A and B from Part A to design a
circuit that will exhibit the same behavior as Figure~\ref{fig:part_c}. Using
additional circuits if needed, an op-amp and resistors from 1\,k$\Omega$ to
10\,k$\Omega$. The following constraints must be satisfied: $V_{OA} = K_a
\cdot V_{IA}$, $V_{OB} = K_b \cdot V_{IB}$, and $V_O = K_a \cdot K_b \cdot K_c
\cdot V_S$. $K_c = 8$.

\section*{Equipment}

\begin{itemize}
    \item Lab Equipment
    \item LTspice
    \item Hantek 3-in-1 Digital Equipment
    \item Powered Breadboard
    \item 1\,k, 3\,k, 4\,k, 4.3\,k, 6\,k $\Omega$ resistors
    \item Operational Amplifier
\end{itemize}

\section*{Simulations / Hand Calculations}

\subsection*{Part A}

For Part A, we needed to design a circuit that could satisfy $V_{OA} = K_a
\cdot V_{IA}$, and another circuit that satisfies $V_{OB} = K_b \cdot V_{IB}$,
with values $K_a = 0.75$ and $K_b = 0.6$.

\begin{figure}[H]
    \centering
    \includegraphics[width=0.85\textwidth]{figures/fig04a_circuit_a_schem.png}
    \caption{Circuit (a) in Part A}
\end{figure}

\begin{figure}[H]
    \centering
    \includegraphics[width=0.6\textwidth]{figures/fig04b_voa_opdata.png}
    \caption{$V_{OA}(0.75\text{V}) = V_{in}(1\text{V}) \cdot K_a(0.75)$}
\end{figure}

\begin{figure}[H]
    \centering
    \includegraphics[width=0.75\textwidth]{figures/fig04c_waveform_a.png}
    \caption{Waveform for square-wave voltage}
\end{figure}

For Circuit A, we determined the resistor values 1\,k$\Omega$ and 3\,k$\Omega$
would satisfy the requirements.

\begin{figure}[H]
    \centering
    \includegraphics[width=0.85\textwidth]{figures/fig05a_circuit_b_schem.png}
    \caption{Circuit (b) in Part A}
\end{figure}

\begin{figure}[H]
    \centering
    \includegraphics[width=0.6\textwidth]{figures/fig05b_vob_opdata.png}
    \caption{$V_{OB}(0.6\text{V}) = K_b(0.6) \cdot V_{in}(1\text{V})$}
\end{figure}

\begin{figure}[H]
    \centering
    \includegraphics[width=0.75\textwidth]{figures/fig05c_waveform_b.png}
    \caption{Waveform for square-wave voltage}
\end{figure}

For Circuit B, we determined that 4\,k$\Omega$ and 6\,k$\Omega$ were good
resistor values to yield the proper results.

\begin{figure}[H]
    \centering
    \includegraphics[width=0.7\textwidth]{figures/fig06_handcalc_parta.png}
    \caption{Hand calculations for both circuits in Part A}
\end{figure}

\subsection*{Part B}

\begin{figure}[H]
    \centering
    \includegraphics[width=0.9\textwidth]{figures/fig07a_cascade_schem.png}
    \caption{Cascaded original circuit outputs 418\,mV, so $V_{OB} = K_a
    \cdot K_b \cdot V_{IA}$ does not apply}
\end{figure}

\begin{figure}[H]
    \centering
    \includegraphics[width=0.6\textwidth]{figures/fig07b_voutb_opdata.png}
    \caption{$V_{OUT(B)} = 418.6$\,mV}
\end{figure}

\begin{figure}[H]
    \centering
    \includegraphics[width=0.75\textwidth]{figures/fig07c_remodeled_schem.png}
    \caption{Remodeled version of Circuit B (3\,k$\Omega$ changed to
    4.3\,k$\Omega$)}
\end{figure}

After concluding the original circuit resulted in a voltage output that
didn't align with the equation, we modified circuit (a) and changed the
3\,k$\Omega$ resistor to a 4.3\,k$\Omega$ one. As shown below, it reflects
450\,mV, which is correct for the given constraint.

\begin{figure}[H]
    \centering
    \includegraphics[width=0.6\textwidth]{figures/fig07d_vob450_opdata.png}
    \caption{$V_{OB}(450\text{mV}) = V_{IA}(1\text{V}) \cdot K_a(0.75) \cdot
    K_b(0.6)$}
\end{figure}

\begin{figure}[H]
    \centering
    \includegraphics[width=0.75\textwidth]{figures/fig07e_waveform_450.png}
    \caption{Voltage waveform for square wave}
\end{figure}

\begin{figure}[H]
    \centering
    \includegraphics[width=0.7\textwidth]{figures/fig08_handcalc_partb.png}
    \caption{Hand calculations for Part B}
\end{figure}

Since $V_{OB} = V_{OA}$, is $V_{OB} = K_a \cdot K_b \cdot V_{IA}$? Checking:
$418\text{mV} \overset{?}{=} (0.75)(0.6)(1\text{V}) = 450$\,mV. No, because the
new Circuit B acts as a load, so the overall resistance in the circuit
increases. Redesigning to fit the constraints (changing the 3\,k$\Omega$
resistor to 4.3\,k$\Omega$) confirms $V_{OA} = 0.75$\,V (750\,mV) and $V_{OB}
= 0.45$\,V (450\,mV), matching $V_{OB} = K_a \cdot K_b \cdot V_{IA}$.

\subsection*{Part C}

\begin{figure}[H]
    \centering
    \includegraphics[width=0.95\textwidth]{figures/fig09_partc_schem_opdata.png}
    \caption{Part C op-amp circuit with simulated operating-point data}
\end{figure}

For Part C of the experiment, we place an op-amp between Circuits A and B. We
then design the op-amp to have a gain of 8. For this, we use a noninverting
configuration. To find the values for the feedback resistor and the resistor
connected to ground, we use the relationship $V_{out} = V_{in}(1 + R_1/R_2)$.
We found that $R_2$ needs to be 7 times greater than $R_1$. To do this, we
utilize the values 7\,k$\Omega$ and 1\,k$\Omega$ as a result.

\section*{Experiment}

\subsection*{Part A}

Measurement of $V_{OA}$ of circuit A $= 0.75698$\,V:

\begin{figure}[H]
    \centering
    \includegraphics[width=0.75\textwidth]{figures/fig10_breadboard_circuita.png}
    \caption{Breadboard implementation of Circuit A}
\end{figure}

\begin{figure}[H]
    \centering
    \includegraphics[width=0.8\textwidth]{figures/fig11_circuita_opdata.png}
    \caption{Circuit A schematic with simulated operating-point data}
\end{figure}

\begin{figure}[H]
    \centering
    \includegraphics[width=0.5\textwidth]{figures/fig12_scope_voa.png}
    \caption{Measured $V_{OA} = 0.75698$\,V}
\end{figure}

Measurement of $V_{OB}$ of circuit B $= 0.60021$\,V:

\begin{figure}[H]
    \centering
    \includegraphics[width=0.75\textwidth]{figures/fig13_breadboard_circuitb.png}
    \caption{Breadboard implementation of Circuit B}
\end{figure}

\begin{figure}[H]
    \centering
    \includegraphics[width=0.8\textwidth]{figures/fig14_circuitb_opdata.png}
    \caption{Circuit B schematic with simulated operating-point data}
\end{figure}

\begin{figure}[H]
    \centering
    \includegraphics[width=0.5\textwidth]{figures/fig15_scope_vob.png}
    \caption{Measured $V_{OB} = 0.60021$\,V}
\end{figure}

The node of $V_{OUT(B)}$ of the noninverting op-amp:

\begin{figure}[H]
    \centering
    \includegraphics[width=0.85\textwidth]{figures/fig16_breadboard_partc.png}
    \caption{Breadboard implementation of the full Part C cascaded circuit
    with op-amp}
\end{figure}

\begin{figure}[H]
    \centering
    \includegraphics[width=0.95\textwidth]{figures/fig17_partc_schem_opdata_repeat.png}
    \caption{Part C op-amp circuit with simulated operating-point data
    (repeated for reference)}
\end{figure}

\begin{figure}[H]
    \centering
    \includegraphics[width=0.5\textwidth]{figures/fig17_scope_gain8.png}
    \caption{Measured output of the gain-of-8 op-amp stage, $3.65845$\,V}
\end{figure}

\section*{Results}

Across the different parts of the experiment, we observe that the measured
results closely match the simulated values, confirming the designs to be
true. In Part A, the measured output of 0.757\,V is closely accurate with the
simulated value of 0.75\,V. In Part B, the measured output of 0.60\,V is
accurate with the simulation result. Finally, in Part C, the op-amp circuit
produced an output of 3.658\,V, which is in range of the simulated output of
3.6\,V. Overall, the experimental and simulated results verify the
theoretical designs with negligible errors.

\section*{Conclusion}

This experiment has shown how unrelated circuit stages can be combined to
obtain a certain overall gain. In Part A and Part B, we designed voltage
divider circuits with gains of $K_a = 0.75$ and $K_b = 0.6$. We then cascaded
the circuits to observe how the individual stages would affect the total
gain. Noticing the differences in the gain at certain points, we determined
that we must do a resistor redesign to regain the gains we originally had. In
Part C, we designed an op-amp circuit with a gain factor of $K_c = 8$, which
restores ideal behavior while allowing for additional amplification of
voltage. Overall, the experiment reinforced how theoretical gain
relationships combined with careful cascaded redesigns could be achieved
accurately under real-world conditions.

\end{document}